\section{Introduction}
\label{sec:intro}

\IEEEPARstart{I}{nduction} motors underpin industrial drivetrains, and their
unplanned failures propagate directly into downtime and safety risk.
Data-driven condition monitoring has consequently produced a large
literature whose reported accuracies now cluster near ceiling on standard
laboratory benchmarks. Deployment experience is harder to reconcile with
those numbers: a model commissioned at one operating point degrades at
another, and faults in service rarely arrive one at a time. A bearing defect
that develops alongside a winding short circuit, broken rotor bars, or
eccentricity confronts the monitoring system with a \emph{compound} state
whose exact combination was, in general, never present in the training
history.

Conventional closed-set formulations cannot even express this situation: a
previously unseen combination $A{+}B$ constitutes an unseen class. The
compound-fault literature therefore moved toward decoupling formulations
that predict elementary mechanisms rather than combinations, and, more
recently, toward zero-shot settings in which compound states must be
recognized from single-fault training data alone. As
Section~\ref{sec:related} reviews, however, that literature has evaluated
almost exclusively \emph{homogeneous} compounds---bearing--bearing or
bearing--gear pairings whose constituents are all observable in the same
vibration channel---and its reported accuracies of roughly 75--87\% have
been obtained under evaluation practices that a parallel line of work on
data leakage has shown to be fragile: sample-level splits, steady operating
conditions, and metrics reported without reference to what label prevalence
alone would achieve.

This paper asks a deliberately harder question on a public dataset that
makes it answerable. The MCC5-THU motor dataset \cite{chen2026} provides 288
recordings of a 2.2-kW three-phase induction motor across 24 diagnostic
states under twelve speed/load conditions, with synchronized triaxial
vibration, three-phase currents, torque, and a key-phase tachometer. Its
nine compound faults pair a mechanical bearing defect with a second
mechanism---inner or outer race, broken rotor bars, static or dynamic
eccentricity, or a stator winding short circuit---so that, in eight of nine
combinations, the two constituents deposit their primary evidence in
\emph{different} sensing modalities. To our knowledge, neither this
cross-modal compound regime nor this dataset has previously received a
constituent-level evaluation.

The study is designed so that every claim is bounded by an explicit control.
The formulation is constituent-level: nine elementary mechanisms are
detected independently by regularized linear probes operating on curated,
mechanism-specific physical-signature families derived from angle-resampled
order spectra, envelope-order descriptors, and phase-current sequence
measures. A crossed protocol varies compound-context exposure over three
regimes---recombination, one-sided, and two-sided compound-context
zero-shot---and crosses each with seen and unseen operating conditions,
under run-level, severity-aware splits with all preprocessing, calibration,
and threshold selection confined to training data. Every reported metric is
accompanied by input-independent constant-predictor references, so that
prevalence-driven scores cannot masquerade as diagnostic skill; oracle
detectability analyses are paired with recording-sensitivity and same-date
controls, so that separability is not silently attributed to fault
mechanisms; and paired run-clustered bootstrap inference with multiplicity
correction distinguishes supported effects from noise.

The contributions are as follows.
\begin{enumerate}
  \item \textbf{A constituent-level formulation and crossed
  composition--context--condition protocol} for compound motor faults whose
  constituents span mechanical and electrical modalities, separating novel
  composition, compound-context deprivation, and operating-condition
  novelty in one design
  (Sections~\ref{sec:methods}--\ref{sec:results}).
  \item \textbf{Prior-aware evaluation}: all constituent metrics are
  reported against input-independent prevalence references, which expose
  several apparently favorable partial-recovery metrics as at or below what
  a constant predictor achieves, while establishing that constituent recall
  and all-constituent recovery genuinely exceed those floors.
  \item \textbf{A controlled analysis of constituent transfer}: within-pair
  oracle detectability, cross-partner transfer, physical-signature
  preservation, and a statistically supported score-orientation reversal,
  each bounded by recording-sensitivity and same-date controls that keep
  session effects from being read as physics.
  \item \textbf{Sensitivity and robustness analyses} over the
  decision-threshold objective, single-fault severity merging, and four
  classifier families, showing that the qualitative generalization pattern
  is stable while absolute exact-match decomposition is an operating-point
  and model property, not an intrinsic bound of the protocol.
  \item \textbf{Documentation of dataset-release properties that materially
  affect protocol design}---recording-session separability,
  severity-variant leakage paths, and label abbreviations---together with
  reference operating-condition experiments
  (Appendix~\ref{app:condition}) that motivate the crossed design.
\end{enumerate}

The principal findings are that closed-set discriminability (accuracy
0.913) coexists with exact constituent-set recovery of 0.083--0.157 on
unseen compounds for the locked probe; that recovery of at least one
constituent is nearly universal while exact decomposition fails, so the
practical failure mode is incomplete decomposition with over-diagnosis
rather than silence; that compound-context deprivation becomes measurably
harmful specifically under simultaneous operating-condition novelty; and
that constituent evidence is often, but directionally and heterogeneously,
transferable across compound partners.

The remainder of the paper is organized as follows.
Section~\ref{sec:related} reviews related work and positions the study.
Section~\ref{sec:methods} details the constituent representation, feature
construction, detection, protocol, metrics, and statistical analysis.
Section~\ref{sec:results} reports results. Section~\ref{sec:conclusion}
concludes. Appendix~\ref{app:condition} reports the operating-condition
reference experiments.
